\begin{abstract}

    % Francais
    Ce travail de bachelor, en partenariat avec l'entreprise Breitling SA, a porté sur la conception d'un banc de test horloger permettant de mesurer la réserve de marche d'un mouvement de montre mécanique de manière accélérée. 
La méthode classique de mesure, basée sur l'observation du fonctionnement du mouvement jusqu'à son arrêt après un remontage complet, est malheureusement très chronophage. 
En effet, certaines montres de la marque fonctionnent pendant plus de 100 heures. 
L'approche de mesure accélérée, étudiée puis proposée par Breitling pour ce travail, consiste à analyser uniquement les 24 dernières heures de la réserve de marche, et fait ainsi office de solution concrète à la problématique évoquée précédemment.

Le projet a consisté à analyser le besoin réel du client afin de définir les exigences fonctionnelles et techniques du projet, puis à concevoir une solution intégrant un posage adaptable, un système motorisé servant à remonter et à dévider de manière contrôlée le mouvement par sa tige de remontoir, ainsi qu'un dispositif de détection automatisé de l'arrêt du mouvement.

Finalement, la réalisation du prototype fonctionnel et de son logiciel de commande a permis d'aboutir à un travail de référence répondant aux attentes de l'entreprise.  
De plus, ce travail constitue une base solide sur laquelle l'entreprise pourra s'appuyer, qu'elle choisisse de s'en inspirer ou non pour finaliser les derniers ajustements nécessaires au développement de leur propre version industrielle.

%% L'asterisme est un signe typographique en forme d'étoile, utilisé pour marquer une pause dans un texte ou pour séparer des paragraphes. Il est souvent utilisé pour indiquer un changement de scène dans un récit. Bien qu'il se fasse rare dans la typographie moderne, c'est un symbole de choix pour séparer les différentes langues du résumé de thèse.
    \asterism

    % English

\end{abstract}